\chapter{Building a Web Service}\label{ch:building-web-service}

\index{web service}
\index{HTTP server}

In the previous chapter we built a command-line tool.
Now we are going to build something that talks to the world: a web service.
By the end of this chapter we will have a running application with HTML pages, a JSON API, database persistence, and middleware---all written in Lattice.

We will use four pieces of the Lattice ecosystem:
the \lstinline|http_server| library for routing and request handling,
the \lstinline|template| library for rendering HTML,
the \lstinline|orm| library for SQLite persistence,
and the built-in \lstinline|json_parse|/\lstinline|json_stringify| functions for JSON APIs.
Let us start at the front door.

%% ─────────────────────────────────────────────────────────────────────────────
\section{The \texttt{http\_server} Library}\label{sec:http-server-library}
%% ─────────────────────────────────────────────────────────────────────────────

\index{http\_server library}
\index{routing!HTTP}

The \lstinline|http_server| library (in \filename{lib/http\_server.lat}) builds on Lattice's raw TCP networking primitives---\lstinline|tcp_listen|, \lstinline|tcp_accept|, \lstinline|tcp_read|, \lstinline|tcp_write|, and \lstinline|tcp_close|---to provide a high-level, Express-style API for building HTTP applications.

\subsection{Hello, Web}\label{sec:http-hello-web}

The smallest possible web server fits in a handful of lines:

\begin{lstlisting}[language=Lattice, caption={A minimal Lattice web server}]
import "lib/http_server" as http

let app = http.new()

app.get("/", |req, res| {
    return http.response(200, "Hello from Lattice!")
})

app.listen(8080)
\end{lstlisting}

Run this, open a browser to \texttt{http://localhost:8080}, and you will see the greeting.
Let us break down what happens:

\begin{enumerate}
  \item \lstinline|http.new()| creates an application Map containing route tables and a middleware stack.
  \item \lstinline|app.get("/", handler)| registers a handler for GET requests to \texttt{/}.
  \item The handler receives a request Map (\lstinline|req|) and a response context Map (\lstinline|res|), and returns a response Map built by \lstinline|http.response|.
  \item \lstinline|app.listen(8080)| opens a TCP socket and enters a loop that accepts connections, parses HTTP requests, runs the middleware pipeline, dispatches to the matching handler, and sends back the formatted response.
\end{enumerate}

\begin{defnbox}{The Request Map}
\index{request Map}
Every handler receives a \lstinline|req| Map with these keys:
\begin{itemize}
  \item \lstinline|"method"| --- the HTTP method (\lstinline|"GET"|, \lstinline|"POST"|, etc.)
  \item \lstinline|"path"| --- the URL path without query string
  \item \lstinline|"raw_path"| --- the full URL including query string
  \item \lstinline|"query"| --- a Map of parsed query parameters
  \item \lstinline|"headers"| --- a Map of HTTP headers
  \item \lstinline|"body"| --- the raw request body as a string
  \item \lstinline|"cookies"| --- a Map of parsed cookie name-value pairs
  \item \lstinline|"peer"| --- the client's address
\end{itemize}
\end{defnbox}

\subsection{Response Helpers}\label{sec:http-response-helpers}

\index{http\_server!response helpers}

The library provides several functions for building response Maps:

\begin{lstlisting}[language=Lattice, caption={Response helper functions}]
import "lib/http_server" as http

// Plain text (Content-Type: text/plain)
let text_res = http.response(200, "OK")

// JSON (Content-Type: application/json)
let data = Map::new()
data.set("status", "healthy")
data.set("uptime", 3600)
let json_res = http.json(200, data)

// HTML (Content-Type: text/html)
let html_res = http.html(200, "<h1>Welcome</h1>")

// Redirect (302 by default)
let redir = http.redirect("/dashboard")

// Redirect with custom status
let perm_redir = http.redirect("/new-path", 301)
\end{lstlisting}

Each helper returns a Map with \lstinline|"status"|, \lstinline|"body"|, \lstinline|"headers"|, and \lstinline|"_cookies"| keys.
The server serializes this into a proper HTTP response before sending it over the wire.

\subsection{Route Registration}\label{sec:http-routes}

\index{routing!route registration}

The application object provides registration closures for all standard HTTP methods:

\begin{lstlisting}[language=Lattice, caption={Registering routes for different HTTP methods}]
import "lib/http_server" as http

let app = http.new()

app.get("/items", |req, res| {
    return http.json(200, get_all_items())
})

app.post("/items", |req, res| {
    let body = req.get("body")
    let item = json_parse(body)
    return http.json(201, create_item(item))
})

app.put("/items/update", |req, res| {
    let body = req.get("body")
    let item = json_parse(body)
    return http.json(200, update_item(item))
})

app.delete("/items/remove", |req, res| {
    return http.json(200, remove_item(req))
})

// Register for ALL methods at once
app.all("/health", |req, res| {
    return http.response(200, "OK")
})
\end{lstlisting}

The available methods are \lstinline|app.get|, \lstinline|app.post|, \lstinline|app.put|, \lstinline|app.delete|, \lstinline|app.patch|, \lstinline|app.options|, and \lstinline|app.all|.
If a request arrives for a registered path but an unregistered method, the server returns 405 Method Not Allowed.
If the path itself is not found, it returns 404 Not Found.

\subsection{Query Parameters}\label{sec:http-query-params}

\index{query parameters}

Query strings are automatically parsed into the \lstinline|"query"| Map on the request:

\begin{lstlisting}[language=Lattice, caption={Reading query parameters}]
// Request: GET /search?q=lattice&page=2
app.get("/search", |req, res| {
    let query = req.get("query")
    let search_term = query.get("q")      // "lattice"
    let page = to_int(query.get("page"))   // 2
    // ... perform search ...
    return http.json(200, results)
})
\end{lstlisting}

\subsection{Middleware}\label{sec:http-middleware}

\index{middleware}

Middleware functions run before your route handlers, forming a pipeline that can inspect requests, modify responses, or short-circuit the chain entirely.
Register middleware with \lstinline|app.use|:

\begin{lstlisting}[language=Lattice, caption={Using middleware}]
import "lib/http_server" as http

let app = http.new()

// Built-in request logger
app.use(http.logger())

// Built-in CORS middleware
app.use(http.cors())

// Built-in JSON body parser
app.use(http.body_parser())

// Built-in security headers
app.use(http.security_headers())

app.get("/", |req, res| {
    return http.response(200, "Hello!")
})

app.listen(3000)
\end{lstlisting}

\begin{defnbox}{Middleware Signature}
\index{middleware!signature}
A middleware is a function with the signature \lstinline{|req, res, next|}.
To continue the pipeline, call \lstinline|next(req, res)| and return its result.
To short-circuit (for example, to reject an unauthenticated request), return a response directly without calling \lstinline|next|.
\end{defnbox}

Here is a custom middleware that checks for an API key:

\begin{lstlisting}[language=Lattice, caption={Writing custom middleware}]
fn require_api_key() -> Fn {
    return |req, res, next| {
        let headers = req.get("headers")
        flux api_key = ""
        if headers.has("X-API-Key") {
            api_key = headers.get("X-API-Key")
        }
        if api_key != "secret-key-123" {
            return http.json(401, "Unauthorized")
        }
        return next(req, res)
    }
}

app.use(require_api_key())
\end{lstlisting}

The built-in middleware that ships with the library includes:

\begin{itemize}
  \item \lstinline|http.logger()| --- logs method, path, status, and response time.
  \item \lstinline|http.cors()| --- handles CORS preflight requests and adds appropriate headers. Accepts optional configuration via \lstinline|http.cors_opts()|.
  \item \lstinline|http.body_parser()| --- automatically parses JSON request bodies when the \lstinline|Content-Type| is \lstinline|application/json|, storing the result in \lstinline|req.get("parsed_body")|.
  \item \lstinline|http.security_headers()| --- adds \lstinline|X-Content-Type-Options|, \lstinline|X-Frame-Options|, and other security headers.
  \item \lstinline|http.compression()| --- sets \lstinline|Vary: Accept-Encoding| for correct proxy caching.
\end{itemize}

\subsection{Cookies}\label{sec:http-cookies}

\index{cookies}

The library provides helpers for reading and setting cookies:

\begin{lstlisting}[language=Lattice, caption={Cookie handling}]
app.get("/profile", |req, res| {
    // Read cookies from the request
    let cookies = req.get("cookies")
    flux username = "Guest"
    if cookies.has("username") {
        username = cookies.get("username")
    }

    let r = http.html(200, "<h1>Hello, " + username + "</h1>")

    // Set a cookie on the response
    let opts = http.cookie_opts()  // sensible defaults
    opts.set("max_age", 86400)     // 1 day
    return http.set_cookie(r, "visited", "true", opts)
})

app.get("/logout", |req, res| {
    let r = http.redirect("/")
    let opts = http.cookie_opts()
    return http.clear_cookie(r, "username", opts)
})
\end{lstlisting}

The \lstinline|http.cookie_opts()| function returns a Map with sensible defaults: \lstinline|Path=/|, \lstinline|HttpOnly=true|, \lstinline|SameSite=Lax|.
You can override any of these before passing the Map to \lstinline|set_cookie|.

\begin{warnbox}{Maps Are Pass-by-Value}
Since Lattice Maps are pass-by-value, \lstinline|http.set_cookie| returns a \emph{new} response Map.
Always use the return value: \lstinline|let r = http.set_cookie(r, ...)|.
Forgetting this is a common source of ``my cookies are not being set'' bugs.
\end{warnbox}

%% ─────────────────────────────────────────────────────────────────────────────
\section{The \texttt{template} Library for HTML Templating}\label{sec:template-library}
%% ─────────────────────────────────────────────────────────────────────────────

\index{template library}
\index{HTML templating}

Building HTML strings by hand is tedious and error-prone.
The \lstinline|template| library (in \filename{lib/template.lat}) provides a Mustache/Jinja2-inspired template engine with variable interpolation, conditionals, loops, filters, includes, and template inheritance.

\subsection{Variables and Escaping}\label{sec:template-variables}

\index{template!variables}

The core operation is variable interpolation with double-braces:

\begin{lstlisting}[language=Lattice, caption={Template variable interpolation}]
import "lib/template" as tmpl

let t = tmpl.compile("Hello, {{name}}!")

let data = Map::new()
data.set("name", "World")

let result = tmpl.render(t, data)
// result: "Hello, World!"
\end{lstlisting}

By default, \lstinline|{{name}}| HTML-escapes the output---so if \lstinline|name| is \lstinline|"<script>alert('xss')</script>"|, it becomes \lstinline|"&lt;script&gt;..."|.
When you need raw, unescaped output (for example, pre-rendered HTML), use triple-braces:

\begin{lstlisting}[language=Lattice, caption={Raw vs.\ escaped output}]
let t = tmpl.compile("{{{html_content}}}")

let data = Map::new()
data.set("html_content", "<b>Bold</b>")

let result = tmpl.render(t, data)
// result: "<b>Bold</b>"  (not escaped)
\end{lstlisting}

Nested map access uses dot notation:

\begin{lstlisting}[language=Lattice, caption={Dot notation for nested values}]
let t = tmpl.compile("{{user.name}} ({{user.email}})")

let user = Map::new()
user.set("name", "Alice")
user.set("email", "alice@example.com")

let data = Map::new()
data.set("user", user)

let result = tmpl.render(t, data)
// result: "Alice (alice@example.com)"
\end{lstlisting}

\subsection{Filters}\label{sec:template-filters}

\index{template!filters}

Filters transform a value inline using the pipe syntax:

\begin{lstlisting}[language=Lattice, caption={Using template filters}]
let t = tmpl.compile("""
  Name: {{name | upper}}
  Bio: {{bio | truncate(50)}}
  Tags: {{tags | join(", ")}}
  Role: {{role | default("member")}}
""")
\end{lstlisting}

The available filters are:

\begin{center}
\begin{tabular}{lp{7.5cm}}
\toprule
\textbf{Filter} & \textbf{Description} \\
\midrule
\lstinline|upper|           & Converts to uppercase \\
\lstinline|lower|           & Converts to lowercase \\
\lstinline|capitalize|      & Capitalizes the first letter \\
\lstinline|title|           & Title-cases each word \\
\lstinline|trim|            & Strips leading/trailing whitespace \\
\lstinline|escape_html|     & HTML-escapes the value \\
\lstinline|length|          & Returns the length \\
\lstinline|string|          & Converts to a string \\
\lstinline|reverse|         & Reverses the string \\
\lstinline|default("val")|  & Uses \texttt{val} if the variable is nil or empty \\
\lstinline|truncate(n)|     & Truncates to \texttt{n} characters, adding \texttt{...} \\
\lstinline|replace("a","b")| & Replaces occurrences of \texttt{a} with \texttt{b} \\
\lstinline|join(", ")|      & Joins array elements with the given separator \\
\bottomrule
\end{tabular}
\end{center}

Filters can be chained: \lstinline|{{name | trim | upper}}| first trims, then uppercases.

\subsection{Conditionals}\label{sec:template-conditionals}

\index{template!conditionals}

The library supports both Mustache-style and Jinja2-style conditionals:

\begin{lstlisting}[language=Lattice, caption={Mustache-style conditionals}]
let t = tmpl.compile("""
  {{#if logged_in}}
    Welcome back, {{username}}!
  {{else}}
    Please log in.
  {{/if}}
""")
\end{lstlisting}

\begin{lstlisting}[language=Lattice, caption={Jinja2-style conditionals with elif}]
let t = tmpl.compile("""
  {% if role %}
    {% if is_admin %}
      <span class="badge">Admin</span>
    {% else %}
      <span class="badge">User</span>
    {% endif %}
  {% endif %}
""")
\end{lstlisting}

Truthiness follows Lattice conventions: \lstinline|nil|, \lstinline|false|, \lstinline|0|, \lstinline|0.0|, and \lstinline|""| are all falsy.
Everything else is truthy.

The negated conditional \lstinline|{{#unless ...}}| works like \lstinline|{{#if ...}}| with inverted logic:

\begin{lstlisting}[language=Lattice, caption={Using unless}]
let t = tmpl.compile("""
  {{#unless errors}}
    <p>No errors found.</p>
  {{/unless}}
""")
\end{lstlisting}

\subsection{Loops}\label{sec:template-loops}

\index{template!loops}

Iteration comes in two flavors.
Mustache-style \lstinline|each| works well for arrays of maps:

\begin{lstlisting}[language=Lattice, caption={Mustache-style each loop}]
let t = tmpl.compile("""
  <ul>
  {{#each items}}
    <li>{{name}} - {{price}}</li>
  {{/each}}
  </ul>
""")
\end{lstlisting}

Inside an \lstinline|each| block, the keys of each Map item are promoted into scope.
The special variables \lstinline|{{@index}}|, \lstinline|{{@first}}|, and \lstinline|{{@last}}| are also available.

Jinja2-style \lstinline|for| loops give you a named loop variable:

\begin{lstlisting}[language=Lattice, caption={Jinja2-style for loop}]
let t = tmpl.compile("""
  <table>
  {% for user in users %}
    <tr class="{% if loop.first %}first{% endif %}">
      <td>{{loop.index1}}</td>
      <td>{{user.name}}</td>
    </tr>
  {% endfor %}
  </table>
""")
\end{lstlisting}

The \lstinline|loop| variable inside a Jinja2 for-loop provides:
\lstinline|loop.index| (0-based),
\lstinline|loop.index1| (1-based),
\lstinline|loop.first|,
\lstinline|loop.last|, and
\lstinline|loop.length|.

\subsection{Includes and Template Inheritance}\label{sec:template-includes}

\index{template!includes}
\index{template!inheritance}

For larger applications, you can split templates into partials and base layouts:

\begin{lstlisting}[language=Lattice, caption={Including a partial}]
// In your template:
// {% include "templates/header.html" %}
// ... page content ...
// {% include "templates/footer.html" %}
\end{lstlisting}

Template inheritance lets child templates override blocks defined in a parent:

\begin{lstlisting}[language={}, caption={base.html --- a base layout}]
<html>
<head><title>{% block title %}My App{% endblock %}</title></head>
<body>
  <nav>...</nav>
  <main>{% block content %}{% endblock %}</main>
  <footer>...</footer>
</body>
</html>
\end{lstlisting}

\begin{lstlisting}[language={}, caption={page.html --- a child template}]
{% extends "templates/base.html" %}
{% block title %}Dashboard{% endblock %}
{% block content %}
  <h1>Dashboard</h1>
  <p>Welcome back!</p>
{% endblock %}
\end{lstlisting}

When you compile and render \texttt{page.html}, the engine loads \texttt{base.html}, finds the block definitions in the child, and substitutes them.
The result is a complete HTML page with ``Dashboard'' as the title and the dashboard content in the \lstinline|<main>| element.

\begin{lstlisting}[language=Lattice, caption={Rendering an inherited template}]
import "lib/template" as tmpl

let page = tmpl.from_file("templates/page.html")
let output = tmpl.render(page, Map::new())
\end{lstlisting}

\begin{tipbox}{Compile Once, Render Many}
The \lstinline|tmpl.compile| function parses a template string into a segment tree.
If you are rendering the same template with different data (as you would in a web server), compile once at startup and call \lstinline|tmpl.render| on each request.
This avoids re-parsing the template on every request.
\end{tipbox}

%% ─────────────────────────────────────────────────────────────────────────────
\section{The \texttt{orm} Library for SQLite Persistence}\label{sec:orm-library}
%% ─────────────────────────────────────────────────────────────────────────────

\index{orm library}
\index{SQLite}
\index{database}

Most web services need to persist data.
The \lstinline|orm| library (in \filename{lib/orm.lat}) provides a lightweight ORM layer over Lattice's SQLite native extension.
It loads the extension via \lstinline|require_ext("sqlite")| and wraps the raw SQL operations in a model-based API.

\subsection{Connecting and Defining Models}\label{sec:orm-connect}

\begin{lstlisting}[language=Lattice, caption={Setting up a database and model}]
import "lib/orm" as orm

// Connect to a file-based database (or ":memory:" for tests)
let db = orm.connect("app.db")

// Define the schema as a Map of column -> SQL type
let schema = Map::new()
schema.set("id", "INTEGER PRIMARY KEY AUTOINCREMENT")
schema.set("title", "TEXT NOT NULL")
schema.set("body", "TEXT")
schema.set("created_at", "TEXT")

// Create a model bound to the "posts" table
let Post = orm.model(db, "posts", schema)

// Create the table (safe to call multiple times)
let create_table = Post.get("create_table")
create_table(0)
\end{lstlisting}

\begin{notebox}{The Dummy Argument}
You may notice that \lstinline|create_table(0)| passes a dummy argument.
This is because the ORM's closure signatures in \filename{lib/orm.lat} accept a parameter (e.g., \lstinline{|_x|}) even when one is not semantically needed.
Simply pass any value---\lstinline|0|, \lstinline|nil|, or \lstinline|""|---to satisfy the call.
\end{notebox}

\subsection{CRUD Operations}\label{sec:orm-crud}

\index{CRUD operations}

The model Map provides closures for all the standard operations:

\begin{lstlisting}[language=Lattice, caption={Creating records}]
let create = Post.get("create")

let post_data = Map::new()
post_data.set("title", "Hello World")
post_data.set("body", "This is my first post.")
post_data.set("created_at", time_format(time(), "%Y-%m-%d %H:%M:%S"))

let post_id = create(post_data)
print("Created post with id: " + to_string(post_id))
\end{lstlisting}

\begin{lstlisting}[language=Lattice, caption={Reading records}]
// Find by id
let find = Post.get("find")
let post = find(post_id)
print(post.get("title"))  // "Hello World"

// Get all records
let all = Post.get("all")
let posts = all(0)
for p in posts {
    print(p.get("title"))
}

// Query with a WHERE clause
let where_fn = Post.get("where")
let recent = where_fn("created_at > ?", ["2025-01-01"])
\end{lstlisting}

\begin{lstlisting}[language=Lattice, caption={Updating and deleting records}]
// Update
let update = Post.get("update")
let changes = Map::new()
changes.set("title", "Hello World (Updated)")
update(post_id, changes)

// Delete
let delete_fn = Post.get("delete")
delete_fn(post_id)

// Count
let count = Post.get("count")
print("Total posts: " + to_string(count(0)))
\end{lstlisting}

The \lstinline|where| closure accepts a SQL condition string with \lstinline|?| placeholders and an array of parameter values.
The parameters are safely bound by the SQLite extension, so you get protection against SQL injection for free.

\begin{lstlisting}[language=Lattice, caption={Closing the database}]
// Always close when done
orm.close(db)
\end{lstlisting}

\begin{warnbox}{Always Close Your Database}
SQLite databases should be closed when your application shuts down to ensure all writes are flushed.
In a long-running web server you typically keep the connection open for the lifetime of the process, but make sure to call \lstinline|orm.close(db)| if you ever exit gracefully.
\end{warnbox}

%% ─────────────────────────────────────────────────────────────────────────────
\section{JSON APIs with \texttt{json\_parse} and \texttt{json\_stringify}}\label{sec:json-apis}
%% ─────────────────────────────────────────────────────────────────────────────

\index{JSON}
\index{json\_parse}
\index{json\_stringify}
\index{API!JSON}

Lattice's JSON support is implemented as built-in functions backed by a C parser and serializer (in \filename{src/json.c}).
The recursive-descent parser handles the full JSON specification, including Unicode escapes, nested structures, and all JSON types.

\subsection{Parsing JSON}\label{sec:json-parsing}

\lstinline|json_parse| takes a JSON string and returns the corresponding Lattice value---a Map for objects, an Array for arrays, and the appropriate scalar types for strings, numbers, booleans, and null:

\begin{lstlisting}[language=Lattice, caption={Parsing JSON strings}]
let data = json_parse("{\"name\":\"Alice\",\"age\":30}")
// data is a Map
print(data.get("name"))  // "Alice"
print(data.get("age"))   // 30

let items = json_parse("[1, 2, 3]")
// items is an Array
print(items[0])  // 1

let nested = json_parse("""
    {
        "users": [
            {"name": "Alice", "active": true},
            {"name": "Bob", "active": false}
        ],
        "total": 2
    }
""")
let users = nested.get("users")
print(users[0].get("name"))  // "Alice"
\end{lstlisting}

JSON null maps to Lattice \lstinline|nil|.
JSON booleans map to Lattice \lstinline|true| and \lstinline|false|.
JSON integers stay as \lstinline|Int|, and numbers with decimal points or exponents become \lstinline|Float|.

\subsection{Serializing to JSON}\label{sec:json-stringify}

\lstinline|json_stringify| converts a Lattice value to a JSON string:

\begin{lstlisting}[language=Lattice, caption={Serializing Lattice values to JSON}]
let data = Map::new()
data.set("name", "Alice")
data.set("scores", [95, 87, 92])
data.set("active", true)

let json_text = json_stringify(data)
print(json_text)
// {"name":"Alice","scores":[95,87,92],"active":true}
\end{lstlisting}

The serializer handles Maps, Arrays, Tuples, Buffers, strings, numbers, booleans, nil, and Refs (which are dereferenced automatically).
Types that have no JSON representation---closures, channels, ranges, structs---will cause an error.

\begin{tipbox}{Round-Tripping}
\lstinline|json_parse(json_stringify(value))| will faithfully round-trip any value composed of Maps, Arrays, strings, numbers, booleans, and nil.
This makes JSON a convenient serialization format for inter-process communication, configuration files, and API payloads.
\end{tipbox}

\subsection{JSON in HTTP Handlers}\label{sec:json-http-handlers}

Combining JSON with the HTTP server is where things get practical.
Here is a pattern you will use constantly:

\begin{lstlisting}[language=Lattice, caption={A JSON API endpoint}]
import "lib/http_server" as http

let app = http.new()
app.use(http.body_parser())

app.post("/api/echo", |req, res| {
    let parsed = req.get("parsed_body")
    if parsed == nil {
        return http.json(400, "Invalid JSON")
    }
    // Echo back the parsed data with a timestamp
    parsed.set("received_at", time_format(time(), "%Y-%m-%dT%H:%M:%S"))
    return http.json(200, parsed)
})
\end{lstlisting}

The \lstinline|http.body_parser()| middleware automatically parses JSON request bodies and stores the result in \lstinline|req.get("parsed_body")|, saving you from calling \lstinline|json_parse| manually on every handler.

%% ─────────────────────────────────────────────────────────────────────────────
\section{A Complete Web Application from Scratch}\label{sec:complete-web-app}
%% ─────────────────────────────────────────────────────────────────────────────

\index{web application!complete example}

Let us bring everything together and build \textbf{Lattice Notes}---a note-taking web application with HTML pages for humans and a JSON API for machines.

\subsection{Architecture}\label{sec:webapp-architecture}

\begin{lstlisting}[language={}, caption={Project structure}]
lattice-notes/
  app.lat                 # main entry point
  templates/
    base.html             # base layout
    index.html            # note listing page
    note.html             # single note page
\end{lstlisting}

The application will have these routes:

\begin{center}
\begin{tabular}{llp{6cm}}
\toprule
\textbf{Method} & \textbf{Path} & \textbf{Description} \\
\midrule
GET  & \texttt{/}               & List all notes (HTML) \\
GET  & \texttt{/notes/view}     & View a single note (HTML, \texttt{?id=N}) \\
POST & \texttt{/notes/create}   & Create a note from form data \\
GET  & \texttt{/api/notes}      & List all notes (JSON) \\
POST & \texttt{/api/notes}      & Create a note (JSON) \\
GET  & \texttt{/api/notes/one}  & Get a note by id (JSON, \texttt{?id=N}) \\
\bottomrule
\end{tabular}
\end{center}

\subsection{Templates}\label{sec:webapp-templates}

First, our base layout:

\begin{lstlisting}[language={}, caption={templates/base.html}]
<!DOCTYPE html>
<html>
<head>
  <meta charset="utf-8">
  <title>{% block title %}Lattice Notes{% endblock %}</title>
  <style>
    body { font-family: sans-serif; max-width: 800px;
           margin: 0 auto; padding: 20px; }
    .note { border: 1px solid #ddd; padding: 12px;
            margin: 8px 0; border-radius: 4px; }
    .note h3 { margin: 0 0 8px 0; }
    form input, form textarea { display: block;
            width: 100%; margin: 8px 0; padding: 8px; }
    form button { padding: 8px 16px; }
  </style>
</head>
<body>
  <h1><a href="/">Lattice Notes</a></h1>
  {% block content %}{% endblock %}
</body>
</html>
\end{lstlisting}

The index page lists all notes and includes a creation form:

\begin{lstlisting}[language={}, caption={templates/index.html}]
{% extends "templates/base.html" %}
{% block title %}All Notes{% endblock %}
{% block content %}
  <h2>All Notes ({{note_count}})</h2>

  {{#each notes}}
    <div class="note">
      <h3><a href="/notes/view?id={{id}}">{{title}}</a></h3>
      <p>{{body | truncate(120)}}</p>
      <small>{{created_at}}</small>
    </div>
  {{/each}}

  {{#unless notes}}
    <p>No notes yet. Create one below!</p>
  {{/unless}}

  <hr>
  <h2>New Note</h2>
  <form method="POST" action="/notes/create">
    <input name="title" placeholder="Title" required>
    <textarea name="body" placeholder="Body" rows="4"></textarea>
    <button type="submit">Create Note</button>
  </form>
{% endblock %}
\end{lstlisting}

And the single-note view:

\begin{lstlisting}[language={}, caption={templates/note.html}]
{% extends "templates/base.html" %}
{% block title %}{{title}}{% endblock %}
{% block content %}
  <h2>{{title}}</h2>
  <p>{{body}}</p>
  <small>Created: {{created_at}}</small>
  <p><a href="/">Back to all notes</a></p>
{% endblock %}
\end{lstlisting}

\subsection{The Application Source}\label{sec:webapp-source}

\begin{lstlisting}[language=Lattice, caption={app.lat --- Lattice Notes web application}]
import "lib/http_server" as http
import "lib/template" as tmpl
import "lib/orm" as orm
import "lib/dotenv" as dotenv
import "lib/log" as log

// ── Configuration ───────────────────────────────────────
dotenv.load()

let log_cfg = Map::new()
log_cfg.set("level", "info")
let logger = log.new(log_cfg)

// ── Database setup ──────────────────────────────────────
let db = orm.connect("notes.db")

let schema = Map::new()
schema.set("id", "INTEGER PRIMARY KEY AUTOINCREMENT")
schema.set("title", "TEXT NOT NULL")
schema.set("body", "TEXT")
schema.set("created_at", "TEXT")

let Note = orm.model(db, "notes", schema)
let create_table = Note.get("create_table")
create_table(0)

logger.info("Database initialized")

// ── Pre-compile templates ───────────────────────────────
let index_tmpl = tmpl.from_file("templates/index.html")
let note_tmpl = tmpl.from_file("templates/note.html")

// ── Helper: parse form body ─────────────────────────────
fn parse_form(body: String) -> Map {
    let result = Map::new()
    let pairs = body.split("&")
    for pair in pairs {
        let eq = pair.index_of("=")
        if eq > 0 {
            let key = pair.substring(0, eq)
            let val = pair.substring(eq + 1, len(pair))
            // Basic URL decoding: + -> space
            let decoded = val.replace("+", " ")
            result.set(key, decoded)
        }
    }
    return result
}

// ── Application ─────────────────────────────────────────
let app = http.new()

// Middleware
app.use(http.logger())
app.use(http.body_parser())
app.use(http.cors())

// ── HTML Routes ─────────────────────────────────────────

// Home page: list all notes
app.get("/", |req, res| {
    let all = Note.get("all")
    let notes = all(0)
    let count = Note.get("count")

    let data = Map::new()
    data.set("notes", notes)
    data.set("note_count", count(0))

    let page = tmpl.render(index_tmpl, data)
    return http.html(200, page)
})

// View a single note
app.get("/notes/view", |req, res| {
    let query = req.get("query")
    if !query.has("id") {
        return http.redirect("/")
    }
    let note_id = to_int(query.get("id"))
    let find = Note.get("find")
    let note = find(note_id)
    if note == nil {
        return http.html(404, "<h1>Note not found</h1>")
    }
    let page = tmpl.render(note_tmpl, note)
    return http.html(200, page)
})

// Create a note from an HTML form
app.post("/notes/create", |req, res| {
    let body = req.get("body")
    let form = parse_form(body)
    let create = Note.get("create")

    let note_data = Map::new()
    note_data.set("title", form.get("title"))
    note_data.set("body", form.get("body"))
    note_data.set("created_at",
        time_format(time(), "%Y-%m-%d %H:%M:%S"))

    create(note_data)
    logger.info("Note created: " + form.get("title"))
    return http.redirect("/")
})

// ── JSON API Routes ─────────────────────────────────────

// List all notes
app.get("/api/notes", |req, res| {
    let all = Note.get("all")
    let notes = all(0)
    return http.json(200, notes)
})

// Get a single note by id
app.get("/api/notes/one", |req, res| {
    let query = req.get("query")
    if !query.has("id") {
        let err = Map::new()
        err.set("error", "Missing 'id' parameter")
        return http.json(400, err)
    }
    let note_id = to_int(query.get("id"))
    let find = Note.get("find")
    let note = find(note_id)
    if note == nil {
        let err = Map::new()
        err.set("error", "Note not found")
        return http.json(404, err)
    }
    return http.json(200, note)
})

// Create a note via JSON
app.post("/api/notes", |req, res| {
    let parsed = req.get("parsed_body")
    if parsed == nil {
        let err = Map::new()
        err.set("error", "Invalid JSON body")
        return http.json(400, err)
    }
    let create = Note.get("create")

    let note_data = Map::new()
    note_data.set("title", parsed.get("title"))
    note_data.set("body", parsed.get("body"))
    note_data.set("created_at",
        time_format(time(), "%Y-%m-%d %H:%M:%S"))

    let new_id = create(note_data)

    let result = Map::new()
    result.set("id", new_id)
    result.set("message", "Note created")
    logger.info("Note created via API: " + parsed.get("title"))
    return http.json(201, result)
})

// ── Start the server ────────────────────────────────────
logger.info("Starting Lattice Notes")
app.listen(8080)
\end{lstlisting}

\subsection{Testing It Out}\label{sec:webapp-testing}

Start the server:

\begin{lstlisting}[language={}, caption={Starting the application}]
$ lattice app.lat
Lattice HTTP server listening on http://localhost:8080
\end{lstlisting}

Open a browser to \texttt{http://localhost:8080} to see the HTML interface.
Create a few notes using the form, then try the JSON API:

\begin{lstlisting}[language={}, caption={Testing the JSON API with curl}]
# List all notes
$ curl http://localhost:8080/api/notes
[{"id":1,"title":"Hello","body":"First note","created_at":"2025-01-15 10:30:00"}]

# Create a note
$ curl -X POST http://localhost:8080/api/notes \
  -H "Content-Type: application/json" \
  -d '{"title":"From the API","body":"Created via curl"}'
{"id":2,"message":"Note created"}

# Get a single note
$ curl http://localhost:8080/api/notes/one?id=1
{"id":1,"title":"Hello","body":"First note","created_at":"2025-01-15 10:30:00"}
\end{lstlisting}

\subsection{Design Walkthrough}\label{sec:webapp-design}

Let us highlight some of the patterns at work:

\begin{enumerate}
  \item \textbf{Templates are compiled once.} We call \lstinline|tmpl.from_file| at startup and reuse the compiled templates on every request.
  This avoids re-parsing the template strings on each page load.

  \item \textbf{Dual interface.} The same data model serves both HTML pages and a JSON API.
  The HTML routes render templates; the API routes call \lstinline|http.json|.
  This dual-interface pattern is common in real applications.

  \item \textbf{Middleware stacking.} The \lstinline|http.logger()| middleware logs every request, \lstinline|http.body_parser()| handles JSON parsing, and \lstinline|http.cors()| adds CORS headers.
  Adding authentication would be one more \lstinline|app.use| call.

  \item \textbf{The ORM hides SQL.} The route handlers never write SQL directly.
  They interact with the \lstinline|Note| model through its \lstinline|create|, \lstinline|find|, \lstinline|all|, and \lstinline|where| closures.
  The \lstinline|where| closure's parameterized queries prevent SQL injection.

  \item \textbf{Error responses are intentional.} Every API endpoint checks for missing parameters and returns structured JSON error objects with appropriate HTTP status codes.
  This makes the API client-friendly.
\end{enumerate}

\begin{notebox}{Streaming and SSE}
The \lstinline|http_server| library also supports chunked transfer encoding via \lstinline|http.stream(status, chunks)| and Server-Sent Events via \lstinline|http.sse_stream(events)|.
These are useful for long-running responses, real-time updates, and progressive data delivery.
Explore \filename{lib/http\_server.lat} for the full API.
\end{notebox}

%% ─────────────────────────────────────────────────────────────────────────────
\section*{Exercises}\label{sec:web-exercises}
%% ─────────────────────────────────────────────────────────────────────────────

\begin{enumerate}
  \item \textbf{Add a DELETE endpoint.}
  Add \lstinline|DELETE /api/notes?id=N| and \lstinline|POST /notes/delete?id=N| routes to the Lattice Notes application.
  The HTML route should redirect back to \texttt{/} after deletion.

  \item \textbf{Note editing.}
  Add a \lstinline|GET /notes/edit?id=N| route that renders an edit form (pre-filled with the existing note data) and a \lstinline|POST /notes/update| route that updates the note and redirects to the view page.

  \item \textbf{Search.}
  Add a \lstinline|GET /search?q=term| route that uses the ORM's \lstinline|where| closure to find notes whose title or body contains the search term.
  Render the results using the index template.

  \item \textbf{Authentication middleware.}
  Write a middleware that checks for a session cookie.
  If the cookie is missing, redirect HTML requests to a \lstinline|/login| page and return 401 for API requests.
  Use \lstinline|http.set_cookie| to issue the session cookie after login.

  \item \textbf{Rate limiting.}
  Build a middleware that tracks the number of requests per client IP address (stored in a Map) and returns \lstinline|429 Too Many Requests| if a client exceeds 60 requests per minute.
  (Hint: store timestamps and use \lstinline|time()| to calculate the window.)
\end{enumerate}

%% ─────────────────────────────────────────────────────────────────────────────
\section*{What's Next}\label{sec:web-whats-next}
%% ─────────────────────────────────────────────────────────────────────────────

Congratulations---you have built both a CLI tool and a web service in Lattice, using real libraries that handle real-world concerns.
You have seen how the language's Map-based closure patterns give libraries like \lstinline|http_server|, \lstinline|template|, and \lstinline|orm| a flexible, composable feel without requiring complex type hierarchies.

These two chapters are the capstone of \emph{The Lattice Handbook}.
Everything we have covered---from the phase system and structured concurrency to modules, testing, and native extensions---comes together when you build real applications.
The appendices that follow serve as a reference for the language grammar, opcode tables, and built-in functions.
Keep them close while you build your next project.
Happy forging.
